\begin{center}
\begin{tikzpicture}[x=1cm,y=1cm]
\draw[rounded corners=2pt,fill=blue!3,draw=blue!22] (0,0) rectangle (6.35,2.2);
\node[anchor=west,font=\small\bfseries,text=newsblue] at (.22,1.90) {铀茂／铀茂阳离子：同一配体骨架的氧化};
\node[font=\large] at (1.38,1.15) {$\mathrm{[U^{IV}(C_8H_8)_2]}$};
\draw[->,thick,newsblue] (2.75,1.18) -- (3.68,1.18);
\node[font=\footnotesize] at (3.22,1.52) {$-\,e^{-}$};
\node[font=\large] at (4.88,1.15) {$\mathrm{[U^{V}(C_8H_8)_2]^{+}}$};
\node[font=\scriptsize,text=newsgray] at (3.2,.48) {两侧为示意化化学式，不表示实验的完整反应条件。};
\end{tikzpicture}
\par\footnotesize\color{newsgray}图 3｜编辑部原创氧化态示意；两个环状配体均为环辛四烯二负离子。来源：Deakin 等，\emph{Nature Communications}，2026-10-23。\color{black}
\end{center}